\hypertarget{estimating-the-threshold-from-a-pilot-budget}{%
\section{\texorpdfstring{Estimating the threshold from a pilot budget}{Estimating the threshold from a pilot budget}}\label{sec:pilot}}

The threshold is exact but its inputs are the state, and an experimenter has snapshots. This
section closes that gap: \(M^*\) can be estimated from the snapshots themselves, without ever
forming the Pauli spectrum whose high-weight part Section~\ref{sec:evaluation-cost} shows is both
necessary to \(\zeta_1\) and badly estimated by shadows. The route around that is to estimate the
two variances \emph{as variances} rather than through their spectral representations.

\hypertarget{four-blocks-and-two-unbiased-estimates}{%
\subsection{Four blocks and two unbiased estimates}\label{four-blocks-and-two-unbiased-estimates}}

Split a pilot of \(M\) snapshots into four disjoint blocks \(A, B, C, D\) and write
\(K_{ij} = \mathrm{Tr}(G_i G_j) = \prod_q \mathrm{Tr}(G_i^q G_j^q)\). Then

\begin{itemize}
\tightlist
\item \(\zeta_2 = \mathrm{Var}[K_{ij}]\) is estimated by the sample variance of \(K\) over the
disjoint pairs \((A_i, B_i)\). Each pair is independent, so this is unbiased directly.
\item \(\mathbb{E}[\mathrm{Tr}(G\rho)^2]\) is estimated by
\(\mathrm{mean}_{i \in A}\, \bar K_{iB}\,\bar K_{iC}\) with
\(\bar K_{iB} = \mathrm{mean}_{j\in B} K_{ij}\). Because \(B\) and \(C\) are disjoint from each
other and from \(A\), the two inner means are independent unbiased estimates of
\(\mathrm{Tr}(G_i\rho)\), so their product is unbiased for \(\mathrm{Tr}(G_i\rho)^2\).
\item \(\mathrm{Tr}(\rho^2)^2\) is estimated by \(\bar K_{AB}\,\bar K_{CD}\) over the four
disjoint blocks, unbiased because the factors are independent and each is unbiased for
\(\mathrm{Tr}(\rho^2)\).
\item \(\zeta_1\) is the difference of the previous two.
\end{itemize}

Nothing here is a plug-in approximation; each piece is exactly unbiased, which is what lets a
small pilot be believed. Every block mean the construction needs is a mean of the
\(4^n\)-dimensional vector \(x_i = \bigotimes_q a_i^q\) of per-qubit Pauli coefficients, since
\(\mathrm{Tr}(G_iG_j) = 2^n \langle x_i, x_j\rangle\) for Hermitian factors, so the whole
computation is \(O(M 4^n)\) with no pairwise matrix ever formed.

The structure also says which factor will be hard. \(\zeta_1\) is a difference of two estimates
sharing the scale \(\mathrm{Tr}(\rho^2)^2\), each carrying noise of order \(\zeta_2/M\); since
\(\zeta_2\) grows as \(7^n\) while \(\zeta_1\) grows as \(\beta^n\) with \(\beta = 5/4\), the
difference is taken between quantities whose common part dominates. That is what the measurements
show: at \(n = 8\) and a pilot of \(32{,}000\) the relative error of \(\hat\zeta_2\) is
\(12.7\%\) while that of \(\hat\zeta_1\) is \(205\%\), and \(\hat\zeta_1\) comes out non-positive
in \(17\) of \(48\) draws, in which case the ratio is undefined rather than large and we report it
as such.

\hypertarget{how-much-pilot-it-takes}{%
\subsection{How much pilot it takes}\label{how-much-pilot-it-takes}}

Table~\ref{tab:pilot} gives the median relative error of \(\widehat{M^*}\) against the exact
statewise \(M^*\) for the same state, on the noisy-pure family. The error falls close to
\(M^{-1/2}\), as an unbiased variance estimate should, and the budget needed for a given accuracy
grows with \(n\) --- but much more slowly than the threshold it is locating.

\begin{table}[htbp]
\centering
\setlength{\tabcolsep}{9pt}
\begin{tabular}{@{}rrrrr@{}}
\toprule
\(n\) & \(M^*\) & \(10\%\) budget & ratio & error at that budget \\
\midrule
\(2\) & \(26\)        & \(8{,}000\)   & \(307\)   & \(5.5\%\) \\
\(3\) & \(110\)       & \(8{,}000\)   & \(72.6\)  & \(6.5\%\) \\
\(4\) & \(541\)       & \(8{,}000\)   & \(14.8\)  & \(8.8\%\) \\
\(5\) & \(2{,}884\)   & \(32{,}000\)  & \(11.1\)  & \(5.8\%\) \\
\(6\) & \(16{,}519\)  & \(32{,}000\)  & \(1.94\)  & \(9.0\%\) \\
\(7\) & \(91{,}741\)  & \(128{,}000\) & \(1.40\)  & \(8.4\%\) \\
\(8\) & \(520{,}955\) & \(512{,}000\) & \(0.98\)  & \(5.7\%\) \\
\bottomrule
\end{tabular}
\caption{Pilot cost against the threshold it locates, noisy-pure ensemble. The \(10\%\) budget is the smallest \emph{swept} budget whose median relative error on \(\widehat{M^*}\) fell below \(10\%\), so it is an upper bound set by the granularity of the grid and not an interpolation; the last column gives the error actually achieved there. Full error-versus-budget curves are in Fig.~\ref{fig:pilotratio}. The ratio crosses unity between \(n = 7\) and \(n = 8\).}
\label{tab:pilot}
\end{table}

The comparison that matters is against \(M^*\) itself, since a pilot larger than the threshold it
locates would be self-defeating. Measured, the ratio of the \(10\%\) budget to \(M^*\) falls from
\(\approx\!300\) at \(n = 2\) through \(1.94\) at \(n = 6\) and \(1.40\) at \(n = 7\) to
\(0.98\) at \(n = 8\): the pilot budget grows about twofold per qubit against the threshold's
roughly fivefold, so the overhead becomes relatively cheaper as \(n\) grows, and it crosses unity
between \(n = 7\) and \(n = 8\). Interpolating log-linearly between those two bracketing measured
points puts the crossing just below \(n = 8\), at \(n \approx 7.95\)
(Fig.~\ref{fig:pilotratio}). We state this as
measured rather than extrapolated because an earlier fit over \(n \le 6\) alone put it near
\(n = 7\) while understating the required budget by a factor of \(1.85\) at \(n = 7\) and
\(3.8\) at \(n = 8\); the crossing survives that correction only because \(M^*\) grows
faster than the pilot does.

The tail of that convergence was worth a second look, because it is where an unbiased estimator
can still mislead. At sixteen repetitions the last octave at \(n = 8\) appeared to flatten,
falling only from \(8.4\%\) to \(7.3\%\), which is what an error floor would look like. It is not
one. Re-running the three largest budgets on independent states at forty repetitions --- \(120\)
pooled draws each --- gives \(27.7\%\), \(10.9\%\) and \(5.7\%\), a fitted exponent of \(-1.14\)
with a \(95\%\) confidence interval of \([-1.38, -0.88]\). That is at least as steep as the
\(M^{-1/2}\) an unbiased variance estimate should show, and it excludes the flattening: the
earlier appearance was an artifact of the repetition count. The median of \(\widehat{M^*}\)
relative to \(M^*\) over the same three budgets is \(1.059\), \(0.963\) and \(1.009\), scattering
about unity rather than settling at an offset, which is the other signature a floor would leave
and does not.

The higher-precision run does move one number, and we report the conservative reading. Its median
deviation at \(256{,}000\) is \(10.9\%\) rather than the \(8.4\%\) the sixteen-repetition run
gave, so that budget now sits just above the \(10\%\) line and the bracket at \(n = 8\) becomes
\(512{,}000\). Two independent determinations straddle the threshold there, which is precisely
the grid-granularity fragility Table~\ref{tab:pilot} is caveated for; taking the larger moves the
ratio at \(n = 8\) from \(0.49\) to \(0.98\) and the crossing from \(n \approx 7.3\) to
\(n \approx 7.95\). Table~\ref{tab:pilot} and Fig.~\ref{fig:pilotratio} use the larger value. The
qualitative claim is unchanged --- the pilot cost falls below the threshold it locates by eight
qubits --- and we prefer the version that claims less. A bracket read off a doubling grid is a
coarse instrument near the threshold, and should be read as one.

The varying-estimand family is harder in absolute terms, and we report it separately rather than
pooled, because pooling would hide the reason. On variable-rank the \(10\%\) budget is
\(8{,}000\) at \(n = 2\), \(32{,}000\) at \(n = 4\), \(128{,}000\) at \(n = 6\), and more than
\(256{,}000\) at \(n = 7\), where the error is still \(19.7\%\) at that budget. But its thresholds
are correspondingly larger --- \(M^* = 503{,}444\) at \(n = 7\) against \(91{,}741\) for
noisy-pure --- so the two move together: rank-deficient states have smaller \(\zeta_1\), which
raises \(M^*\) and makes \(\zeta_1\) harder to estimate at once. The ratio, not the budget, is the
quantity that transfers.

\begin{figure}[tb]
\centering
\includegraphics[width=0.9\linewidth]{fig8_pilot_convergence}\hfill
\includegraphics[width=0.9\linewidth]{fig9_pilot_over_mstar}
\caption{Estimating the threshold from data. \emph{Left:} median relative error of $\widehat{M^*}$ against the exact statewise $M^*$ versus pilot budget, noisy-pure, one curve per size; the dashed line marks $10\%$ accuracy. \emph{Right:} the $10\%$ budget relative to $M^*$ itself. Circles are the smallest swept budget that achieved $10\%$, so they are upper bounds set by the grid, which is why the first three coincide; squares are a log-log fit through the usable range. The star marks where the measured ratio crosses unity, interpolated between the two measured sizes that bracket it. Below the dashed line the pilot costs less than the threshold it locates.}
\label{fig:pilotratio}
\end{figure}

What this does not provide is a performance guarantee. The estimators are unbiased and their
convergence is measured, but bounding the probability that a pilot puts an experiment on the wrong
side of \(M^*\) needs a concentration bound for \(\hat\zeta_1\), which is a difference of degree-2
and degree-4 statistics in snapshots whose kernel variance grows as \(7^n\); the fourth moment
that a Bernstein-type argument would need is itself a higher-order projection variance we do not
have in closed form. We state the convergence we measured and leave the guarantee open.
